\thispagestyle{plain}			% Supress header
\setlength{\parskip}{0pt plus 1.0pt}
\section*{Abstract}
L'ottimizzazione matematica permette di affrontare numerosi problemi decisionali
in ambito industriale, logistico ed economico, ma tra la formulazione di un
problema e la sua risoluzione si colloca la \emph{modellazione}: la traduzione del
problema in un modello matematico trattabile da un risolutore. Gli strumenti
esistenti per questa attività sono spesso commerciali, legati a un linguaggio di
programmazione ospite o poco orientati all'apprendimento. Nell'ecosistema del
linguaggio Rust, inoltre, nessuno di essi consente di esprimere vincoli logici 
in modelli di programmazione lineare intera mista, linearizzandoli
automaticamente.

Questa tesi presenta ROOC, un linguaggio dichiarativo per la modellazione di
problemi di ottimizzazione lineare e l'ambiente che lo accompagna. ROOC adotta una
notazione vicina a quella matematica e separa la struttura del modello dai dati; a
partire da tale descrizione si comporta come un \emph{compilatore}, traducendo il
modello di alto livello in una forma lineare standard attraverso una sequenza di
trasformazioni ispezionabili.

Il lavoro si concentra sul motore di ROOC. I contributi principali sono la
progettazione del linguaggio, dotato di un sistema di tipi e di costrutti di alto
livello quali iteratori, funzioni e tipi di dato strutturati; la
realizzazione di un compilatore multi-fase che riduce il modello alla forma
standard; l'estensione del solver \texttt{microlp}, ottenuto come \emph{fork} del
crate archiviato \texttt{minilp} e ampliato dal solo simplesso al supporto della
programmazione lineare intera mista; e la distribuzione del motore come libreria
riutilizzabile e come piattaforma web a fini didattici, eseguita interamente nel
browser tramite la compilazione in WebAssembly.

% KEYWORDS (MAXIMUM 10 WORDS)
\vfill
Parole chiave: ottimizzazione, linguaggio di modellazione, compilatori, programmazione lineare, MILP, simplesso, WebAssembly, Rust, fluent API.

\thispagestyle{empty}
\mbox{}
